\chapter{Geltungsbereich und Grundbegriffe}
\label{b15guide:ch:foundations}

Dieser Leitfaden ergänzt Band~15 um eine ausgewählte Übersicht
zentraler Definitionen, Charakterisierungen und unmittelbarer
Folgerungen für \emph{totale Ordnungen}. Er enthält keine Beweise
oder Beweisskizzen. Die Nummern mit Präfix $15H$ gehören
ausschließlich zu dieser Ergänzung.

Aufgenommen werden Aussagen, bei denen die vollständige
Vergleichbarkeit einen wesentlichen Beitrag leistet. Eine Aussage
wird nicht allein deshalb aufgenommen, weil ein allgemeiner
Halbordnungssatz auch auf totale Ordnungen anwendbar ist.
Benötigte Grundbegriffe werden knapp definiert.
Der Band ist eine Auswahl bekannter Hauptresultate, keine
vollständige Aufzählung aller bekannten oder ableitbaren Sätze.

\section{Voraussetzungen und Abgrenzung}

Durchgehend bezeichnet $(A,\leq)$ eine beliebige total geordnete
Menge. $A$ darf leer sein; ungebundene Elementvariablen liegen
in $A$. Als mengentheoretische Grundlage genügt ZF ohne
zusätzliches Auswahlaxiom. Endlichkeit einer betrachteten
Teilmenge wird nur dort vorausgesetzt, wo sie ausdrücklich steht.

Es werden keine Dichtheits-, Vollständigkeits-, Wohlordnungs-
oder Diskretheitsaxiome an $A$ gestellt. Auch algebraische,
topologische oder metrische Zusatzstrukturen werden nicht
eingeführt. Der Hauptband enthält ein Zahlenbeispiel;
der Leitfaden bleibt bei abstrakten totalen Ordnungen.

\begin{center}
\small
\begin{tabularx}{\linewidth}{@{}p{.72\linewidth}X@{}}
\toprule
Eigenständiges Thema außerhalb dieses Leitfadens & Band\\
\midrule
Allgemeine Funktionen; Injektivität, Surjektivität, Bijektivität & 5--8\\
Äquivalenzrelationen und Quotienten & 11\\
Allgemeine Theorie der Halbordnungen & 12\\
Allgemeine Schranken-, Infimums- und Supremumssätze & 13\\
Allgemeiner Kalkül vorhandener Paarinfima und Paarsuprema & 14\\
Wohlordnungen und Auswahlaxiom & 16\\
Allgemeine Theorie endlicher Mengen & 20\\
Halbverbände und Verbände als axiomatische Strukturen & 45\\
\bottomrule
\end{tabularx}
\end{center}

Ordnungsabbildungen, Isomorphismen und ihre Funktionseigenschaften
sind ebenfalls kein Thema. Minimum und Maximum werden als
ordnungseigene Auswahl eines der vorgegebenen Elemente benutzt.
Die Theorie beliebiger Infima und Suprema wird dafür nicht benötigt.

\section{Totale und strikte Ordnung}

\begin{definition}[Totale Ordnung]
\label{b15guide:def:linear}
Eine Relation $\leq\ \subseteq A\times A$ heißt
\emph{totale} oder \emph{lineare Ordnung}, wenn für alle
$x,y,z\in A$ gilt:
\[
\begin{aligned}
 &x\leq x &&\text{(Reflexivität)},\\
 &(x\leq y\land y\leq z)\Rightarrow x\leq z
     &&\text{(Transitivität)},\\
 &(x\leq y\land y\leq x)\Rightarrow x=y
     &&\text{(Antisymmetrie)},\\
 &x\leq y\ \lor\ y\leq x &&\text{(Vergleichbarkeit)}.
\end{aligned}
\]
Wie in Band~15 schreiben wir dafür $\TotOrd{A,\leq}$.
\end{definition}

\begin{definition}[Strikte und umgekehrte Vergleiche]
Wir setzen
\[
 x<y\ \Longleftrightarrow\ (x\leq y\land x\ne y),
 \qquad x\geq y\ \Longleftrightarrow\ y\leq x,
 \qquad x>y\ \Longleftrightarrow\ y<x.
\]
$x\not\leq y$ und $x\not<y$ bezeichnen die jeweiligen
Negationen. Alle Vergleiche beziehen sich auf den angegebenen Träger.
\end{definition}

\section{Wegweiser}

Kapitel~\ref{b15guide:ch:comparison} charakterisiert die
Vergleichbarkeit und sammelt ihre logischen Folgen.
Kapitel~\ref{b15guide:ch:extrema} behandelt Extremalelemente,
Minima und Maxima.
Kapitel~\ref{b15guide:ch:intervals} enthält die für eine
lineare Anordnung besonderen Intervall- und Trennungsresultate sowie den Median.
Kapitel~\ref{b15guide:ch:sections} beschreibt Anfangsabschnitte,
Schnitte und unmittelbare Nachbarn.
Kapitel~\ref{b15guide:ch:constructions} nennt Konstruktionen,
die wieder totale Ordnungen liefern.

Die Charakterisierungssätze dürfen zum Erkennen einer totalen
Ordnung mit einer zunächst nur halbgeordneten Menge beginnen.
Das ist ein Vergleich der Voraussetzungen, keine erneute
Entwicklung der allgemeinen Halbordnungstheorie.
